\documentclass{article}
\PassOptionsToPackage{numbers}{natbib}
\usepackage[sglblindworkshop]{neurips_2026}
\usepackage[utf8]{inputenc}
\usepackage[T1]{fontenc}
\usepackage{hyperref}
\usepackage{url}
\usepackage{booktabs}
\usepackage{tabularx}
\usepackage{amsfonts}
\usepackage{microtype}
\usepackage{xcolor}
\workshoptitle{Galileo Hackathon 2026}
% Venue notice only: preserve the official typography, margins and spacing.
% This is an independent hackathon report, not a NeurIPS/ML4PS submission.
\makeatletter
\renewcommand{\@noticestring}{Galileo Hackathon 2026. Participant research report.}
\makeatother
\title{[A specific scientific finding or method]}
\author{[Team name] \\ {[Member names and affiliations]}
\\ Team ID: [team-id] \quad Task slug: [task-slug]}
\begin{document}
\maketitle
\begin{abstract}
[State the scientific question, the initial idea, the most informative experiment,
and the final finding with its main limitation. Replace all bracketed prompts
with your own evidence-based account.]
\end{abstract}

\section{Scientific goal \& initial hypothesis}
[Name the task and explain the scientific objective beyond its leaderboard metric.
State your initial hypothesis, why it was plausible, and what observation would
support or contradict it. Define unfamiliar terms.]

\section{Research trajectory}
[Describe how evidence changed your decisions. Include real unsuccessful attempts
or negative results where they occurred; do not invent failures or fill a quota.
Summarize the most informative decision points (six or fewer is a suggestion).
Link the complete original logs in Section~6.]

\begin{table}[ht]
\caption{Decision trail. Replace or extend rows with actual runs; identifiers must
match the accompanying evidence.}
\label{tab:trajectory}
\centering
\begin{tabularx}{\textwidth}{@{}lXXXX@{}}
\toprule
Run ID & Why & Change & Result & Next decision \\
\midrule
{[R01]} & [Question] & [Intervention] & [Observation] & [Decision] \\
{[R02]} & [Question] & [Intervention] & [Observation] & [Decision] \\
\bottomrule
\end{tabularx}
\end{table}

\section{Final model \& comparisons}
[Explain the final method, training procedure and selection rationale.
Bind this account to an exact \textbf{submission ID}, model artifact or checksum,
and code version. Specify evaluation split, metric name and direction, applicable
gates and their outcomes. Distinguish measured outcomes from untested claims.]

\begin{table}[ht]
\caption{Comparable measurements. Use the official baseline, your final model,
and a meaningful alternative. Mark unmeasured values as not measured; do not
invent results. Add uncertainty or repeated-run information when available.}
\label{tab:comparison}
\centering
\begin{tabularx}{\textwidth}{@{}lXXXX@{}}
\toprule
Method & Split & Metric & Gates & Cost, if known \\
\midrule
Official baseline & [Split] & [Value] & [Outcome] & [Compute/time] \\
Final model & [Split] & [Value] & [Outcome] & [Compute/time] \\
{[Alternative]} & [Split] & [Value] & [Outcome] & [Compute/time] \\
\bottomrule
\end{tabularx}
\end{table}

\section{Key insights}
[Give at most three insights. For each, state a \textbf{claim}, the
\textbf{evidence} supporting it (run/table/figure), and a \textbf{limitation}.
Label a proposed causal explanation as a hypothesis unless the experiments
establish causality. A well-supported negative result is a valid insight.]

\section{Limitations \& next experiment}
[Identify the most consequential uncertainty, failure mode or restricted scope.
Propose a specific next experiment, the competing explanations it would test,
and how different outcomes would change your conclusion.]

\section{Reproducibility, contributions \& AI use}
[Provide the exact training and inference commands, environment and code version,
seeds where relevant, data access and split usage, and links to model artifacts
and complete original experiment logs. Report unavailable evidence explicitly.
Identify each member's contributions. State AI tools and uses, and how their
outputs were checked. The research judgments and core arguments are your own;
see the AI policy in README.md.]

% OPTIONAL: delete this entire section if it does not apply.
% It is part of the four-page body, not an appendix.
\section*{Optional: Reward hacking discovery}
[Explain the scientific goal versus the metric incentive; the triggering
condition; a minimal reproducible example; evidence that a high score can coexist
with scientific failure; and a proposed fix with possible false positives.
Make the reasoning your own and link supporting artifacts.]

% References may extend beyond the four-page body. Do not add narrative appendices.
\clearpage
\begin{thebibliography}{9}
\bibitem{replace} [Replace with accurate references you actually used, or remove
this placeholder if no references are needed. Verify titles, authors and links.]
\end{thebibliography}
\end{document}
